\documentclass[journal,12pt,twocolumn]{IEEEtran}

\usepackage{setspace}
\usepackage{xcolor}
\usepackage{graphicx}
\usepackage{float}
\singlespacing
\usepackage{siunitx}
\usepackage[cmex10]{amsmath}
\usepackage{mathtools}
\usepackage{hyperref}
\usepackage[margin=0.78in,columnsep=0.25in]{geometry}
\usepackage{xurl}
\usepackage{amsthm}
\usepackage{mathrsfs}
\usepackage{txfonts}
\usepackage{stfloats}
\usepackage{cite}
\usepackage{cases}
\usepackage[caption=false]{subfig}
\usepackage{longtable}
\usepackage{multirow}
\usepackage{enumitem}
\usepackage{listings}
\usepackage{tikz}
\usetikzlibrary{shapes,arrows,positioning}
\usepackage{circuitikz}

% Reduce avoidable line- and page-breaking warnings
\emergencystretch=2em
\raggedbottom

\let\vec\mathbf

\DeclareMathOperator*{\Res}{Res}

\renewcommand\thesection{\arabic{section}}
\renewcommand\thesubsection{\thesection.\arabic{subsection}}
\renewcommand\thesubsubsection{\thesubsection.\arabic{subsubsection}}

\renewcommand\thesectiondis{\arabic{section}}
\renewcommand\thesubsectiondis{\thesectiondis.\arabic{subsection}}
\renewcommand\thesubsubsectiondis{\thesubsectiondis.\arabic{subsubsection}}

\hyphenation{op-tical net-works semi-conduc-tor}

\lstset{
    language=Python,
    frame=single,
    breaklines=true,
    columns=fullflexible
}

% Imaginary-unit notation and solution heading
\newcommand{\jj}{\mathrm{j}}
\providecommand{\solution}{\par\noindent\textit{Solution.}\ }
\newcommand{\maybeimage}[1]{\IfFileExists{#1}{\includegraphics[width=\columnwidth]{#1}}{\fbox{\parbox[c][0.85in][c]{0.88\columnwidth}{\centering Figure file not included in this source package.}}}}

\begin{document}

\theoremstyle{definition}
\newtheorem{theorem}{Theorem}[section]
\newtheorem{problem}{Problem}
\newtheorem{proposition}{Proposition}[section]
\newtheorem{lemma}{Lemma}[section]
\newtheorem{corollary}[theorem]{Corollary}
\newtheorem{example}{Example}[section]
\newtheorem{definition}{Definition}[section]

\newcommand{\BEQA}{\begin{eqnarray}}
\newcommand{\EEQA}{\end{eqnarray}}
\newcommand{\define}{\stackrel{\triangle}{=}}
\newcommand{\myvec}[1]{\ensuremath{\begin{pmatrix}#1\end{pmatrix}}}
\newcommand{\mydet}[1]{\ensuremath{\begin{vmatrix}#1\end{vmatrix}}}

\bibliographystyle{IEEEtran}

\providecommand{\nCr}[2]{\,^{#1}C_{#2}}
\providecommand{\nPr}[2]{\,^{#1}P_{#2}}
\providecommand{\mbf}{\mathbf}
\providecommand{\pr}[1]{\ensuremath{\Pr\left(#1\right)}}
\providecommand{\qfunc}[1]{\ensuremath{Q\left(#1\right)}}
\providecommand{\sbrak}[1]{\ensuremath{{}\left[#1\right]}}
\providecommand{\lsbrak}[1]{\ensuremath{{}\left[#1\right.}}
\providecommand{\rsbrak}[1]{\ensuremath{{}\left.#1\right]}}
\providecommand{\brak}[1]{\ensuremath{\left(#1\right)}}
\providecommand{\lbrak}[1]{\ensuremath{\left(#1\right.}}
\providecommand{\rbrak}[1]{\ensuremath{\left.#1\right)}}
\providecommand{\cbrak}[1]{\ensuremath{\left\{#1\right\}}}
\providecommand{\lcbrak}[1]{\ensuremath{\left\{#1\right.}}
\providecommand{\rcbrak}[1]{\ensuremath{\left.#1\right\}}}

\theoremstyle{remark}
\newtheorem{rem}{Remark}

\newcommand{\sgn}{\mathop{\mathrm{sgn}}}
\newcommand{\rect}{\mathop{\mathrm{rect}}}
\newcommand{\sinc}{\mathop{\mathrm{sinc}}}

\providecommand{\abs}[1]{\left\vert#1\right\vert}
\providecommand{\res}[1]{\Res\displaylimits_{#1}}
\providecommand{\norm}[1]{\lVert#1\rVert}
\providecommand{\mtx}[1]{\mathbf{#1}}
\providecommand{\mean}[1]{E\left[ #1 \right]}
\providecommand{\fourier}{\overset{\mathcal{F}}{\rightleftharpoons}}
\providecommand{\ztrans}{\overset{\mathcal{Z}}{\rightleftharpoons}}
\providecommand{\system}[1]{\overset{\mathcal{#1}}{\longleftrightarrow}}

\providecommand{\dec}[2]{\ensuremath{\overset{#1}{\underset{#2}{\gtrless}}}}

\numberwithin{equation}{section}

\makeatletter
\@addtoreset{figure}{problem}
\makeatother

\let\StandardTheFigure\thefigure
\renewcommand{\thefigure}{\theproblem}

\def\putbox#1#2#3{
    \makebox[0in][l]{
        \makebox[#1][l]{}
        \raisebox{\baselineskip}[0in][0in]{
            \raisebox{#2}[0in][0in]{#3}
        }
    }
}
\def\rightbox#1{\makebox[0in][r]{#1}}
\def\centbox#1{\makebox[0in]{#1}}
\def\topbox#1{\raisebox{-\baselineskip}[0in][0in]{#1}}
\def\midbox#1{\raisebox{-0.5\baselineskip}[0in][0in]{#1}}

\vspace{3cm}

\title{Fourier Series and the Circuit Design of a 5\,V USB Charger}
\author{Nehan Mohammad, Anvidha Reddy}

\maketitle
\tableofcontents

\renewcommand{\thefigure}{\theenumi}
\renewcommand{\thetable}{\theenumi}

\bigskip

\begin{abstract}
This report develops the Fourier-series and Fourier-transform analysis of a
full-wave rectified waveform, then connects that analysis to the practical
stages of a 5 V charger: voltage conversion, rectification, smoothing,
regulation, and the USB output interface. Circuit explanations are expanded
for intuitive understanding, while numerical values are distinguished from
reference examples and unverified assumptions.
\end{abstract}

%------------------------------------------------
\section{Periodic Function}
%------------------------------------------------

Let

\begin{align}
    x(t) &= A_0\abs{\sin\brak{2\pi f_0t}}
    \label{eq:10-orig-diff-def}
\end{align}

\begin{enumerate}[
    label=\thesection.\arabic*,
    ref=\thesection.\theenumi
]

% Q1.1
\item Plot $x(t)$.

\solution The plot of $x(t)$ is shown in Fig.~\ref{fig:xt}.

\begin{figure}[!htp]
    \centering
    \maybeimage{codes/q1_1_plot.png}
    \caption{Plot of $x(t)=A_0\abs{\sin(2\pi f_0t)}$.}
    \label{fig:xt}
\end{figure}

% Q1.2
\item Show that $x(t)$ is periodic and find its period.

\solution Note that,

\begin{align}
    x\brak{t+\frac{1}{2f_0}}
    &= A_0\abs{
    \sin\brak{2\pi f_0\brak{t+\frac{1}{2f_0}}}
    }\\
    &= A_0\abs{\sin\brak{2\pi f_0t+\pi}}\\
    &= A_0\abs{-\sin\brak{2\pi f_0t}}\\
    &= A_0\abs{\sin\brak{2\pi f_0t}}\\
    &= x(t).
\end{align}

Hence, $x(t)$ is periodic with fundamental period

\begin{align}
    \boxed{T_0=\frac{1}{2f_0}}.
\end{align}

\end{enumerate}

%------------------------------------------------
\section{Fourier Series}
%------------------------------------------------

Consider $A_0=12$ and $f_0=50$ for all numerical calculations.

\begin{enumerate}[
    label=\thesection.\arabic*,
    ref=\thesection.\theenumi
]

% Q2.1
\item If

\begin{align}
    x(t)
    &= \sum_{k=-\infty}^{\infty}
    c_ke^{\jj2\pi kf_0t}
    \label{eq:one-Z-complex}
\end{align}

show that

\begin{align}
    c_k
    &= f_0\int_{-\frac{1}{2f_0}}^{\frac{1}{2f_0}}
    x(t)e^{-\jj2\pi kf_0t}\,dt.
    \label{eq:one-Z}
\end{align}

\solution For some $n\in\mathbb{Z}$, multiplying
both sides of \eqref{eq:one-Z-complex} by
$e^{-\jj2\pi nf_0t}$ gives

\begin{align}
    x(t)e^{-\jj2\pi nf_0t}
    =
    \sum_{k=-\infty}^{\infty}
    c_ke^{\jj2\pi(k-n)f_0t}.
\end{align}

Now, integrating both sides over the interval
$\left[-\frac{1}{2f_0},\frac{1}{2f_0}\right]$,

\begin{align}
    &\int_{-\frac{1}{2f_0}}^{\frac{1}{2f_0}}
    x(t)e^{-\jj2\pi nf_0t}\,dt
    \nonumber\\
    &=
    \sum_{k=-\infty}^{\infty}
    c_k
    \int_{-\frac{1}{2f_0}}^{\frac{1}{2f_0}}
    e^{\jj2\pi(k-n)f_0t}\,dt.
\end{align}

The integral on the right can be evaluated using the
orthogonality property of complex exponentials:

\begin{align}
    &\int_{-\frac{1}{2f_0}}^{\frac{1}{2f_0}}
    e^{\jj2\pi(k-n)f_0t}\,dt
    \nonumber\\
    &=
    \begin{cases}
        \dfrac{1}{f_0}, & k=n,\\[4pt]
        0, & k\neq n.
    \end{cases}
\end{align}

Note that the integration interval has length $1/f_0$.
Although the fundamental period of the given signal
$x(t)=A_0|\sin(2\pi f_0t)|$ is $1/(2f_0)$, the chosen
Fourier basis $e^{\jj2\pi kf_0t}$ has period $1/f_0$. Thus, the integration interval covers two fundamental periods of the signal and one period of the Fourier basis.

Consequently, all terms in the summation vanish except
the term for which $k=n$. Therefore,

\begin{align}
    &\int_{-\frac{1}{2f_0}}^{\frac{1}{2f_0}}
    x(t)e^{-\jj2\pi nf_0t}\,dt
    \nonumber\\
    &=\frac{c_n}{f_0}.
\end{align}

Hence,

\begin{align}
    \boxed{
    c_n=f_0\int_{-\frac{1}{2f_0}}^{\frac{1}{2f_0}}
    x(t)e^{-\jj2\pi nf_0t}\,dt
    }.
\end{align}

% Q2.2

\item Find $c_k$ for \eqref{eq:10-orig-diff-def}.

\solution Using \eqref{eq:one-Z},

\begin{align}
    c_n
    &=f_0\int_{-\frac{1}{2f_0}}^{\frac{1}{2f_0}}
    A_0\abs{\sin(2\pi f_0t)}
    e^{-\jj2\pi nf_0t}\,dt.
\end{align}

Using Euler's identity,

\begin{align}
    e^{-\jj\theta}=\cos\theta-\jj\sin\theta,
\end{align}

we obtain

\begin{align}
    c_n
    &=f_0A_0\int_{-\frac{1}{2f_0}}^{\frac{1}{2f_0}}
    \abs{\sin(2\pi f_0t)}
    \cos(2\pi nf_0t)\,dt
    \nonumber\\
    &\quad-\jj f_0A_0
    \int_{-\frac{1}{2f_0}}^{\frac{1}{2f_0}}
    \abs{\sin(2\pi f_0t)}
    \sin(2\pi nf_0t)\,dt.
\end{align}

Since $x(t)$ is even and the sine term is odd, the
imaginary part vanishes. Thus,

\begin{align}
    c_n
    &=2f_0A_0\int_0^{\frac{1}{2f_0}}
    \sin(2\pi f_0t)\cos(2\pi nf_0t)\,dt.
\end{align}

Using the identity

\begin{align}
    2\sin A\cos B=\sin(A+B)+\sin(A-B),
\end{align}

we get

\begin{align}
    c_n
    &=f_0A_0\int_0^{\frac{1}{2f_0}}
    \sin(2\pi(n+1)f_0t)\,dt
    \nonumber\\
    &\quad+
    f_0A_0\int_0^{\frac{1}{2f_0}}
    \sin(2\pi(1-n)f_0t)\,dt.
\end{align}

For $n\neq\pm1$, evaluating the integrals gives

\begin{align}
    c_n
    &=\frac{A_0}{2\pi}
    \left[
    \frac{1-(-1)^{n+1}}{n+1}
    +
    \frac{1-(-1)^{1-n}}{1-n}
    \right]
    \nonumber\\
    &=\frac{A_0(1+(-1)^n)}{\pi(1-n^2)}.
\end{align}

For $n=\pm1$, direct evaluation of the original integral
gives $c_n=0$.

Therefore,

\begin{align}
    \boxed{
    c_n=
    \begin{cases}
        \dfrac{2A_0}{\pi(1-n^2)}, & n\text{ even},\\[6pt]
        0, & n\text{ odd}.
    \end{cases}
    }
    \label{eq:ck-xt}
\end{align}

Here, $n=0$ is included in the even case, giving
$c_0=2A_0/\pi$.

% Q2.3
\item Verify \eqref{eq:one-Z-complex} using Python.

\solution The Python code \texttt{codes/2\_3.py} verifies
\eqref{eq:one-Z-complex} by comparing the original signal
with its Fourier-series reconstruction, as shown in
Fig.~\ref{fig:ver-complex}.

\begin{figure}[!ht]
    \centering
    \maybeimage{codes/2_3.png}
    \caption{Verification of the complex Fourier series.}
    \label{fig:ver-complex}
\end{figure}

% Q2.4
\item Show that

\begin{align}
    x(t)
    &=a_0+\sum_{k=1}^{\infty}
    \left[
    a_k\cos(2\pi kf_0t)
    +b_k\sin(2\pi kf_0t)
    \right]
    \label{eq:one-Z-real}
\end{align}

and obtain the formulae for $a_k$ and $b_k$.

\solution From \eqref{eq:one-Z-complex},

\begin{align}
    x(t)
    &=\sum_{k=-\infty}^{\infty}
    c_ke^{\jj2\pi kf_0t}\\
    &=c_0+\sum_{k=1}^{\infty}
    \left[
    c_ke^{\jj2\pi kf_0t}
    +c_{-k}e^{-\jj2\pi kf_0t}
    \right].
\end{align}

Using Euler's identities,

\begin{align}
    e^{\jj\theta}&=\cos\theta+\jj\sin\theta,\\
    e^{-\jj\theta}&=\cos\theta-\jj\sin\theta,
\end{align}

we obtain

\begin{align}
    x(t)
    &=c_0+\sum_{k=1}^{\infty}
    \Big[
    (c_k+c_{-k})\cos(2\pi kf_0t)
    \nonumber\\
    &\qquad+
    \jj(c_k-c_{-k})\sin(2\pi kf_0t)
    \Big].
\end{align}

Comparing this with \eqref{eq:one-Z-real}, we get

\begin{align}
    a_k&=
    \begin{cases}
        c_0, & k=0,\\
        c_k+c_{-k}, & k>0,
    \end{cases}
    \label{eq:ak}
\end{align}

and

\begin{align}
    b_k&=\jj(c_k-c_{-k}),\qquad k>0.
    \label{eq:bk}
\end{align}

Using the coefficient formula \eqref{eq:one-Z}, the
corresponding integral expressions are

\begin{align}
    a_0
    &=f_0\int_{-\frac{1}{2f_0}}^{\frac{1}{2f_0}}
    x(t)\,dt,
    \nonumber\\
    a_k
    &=2f_0\int_{-\frac{1}{2f_0}}^{\frac{1}{2f_0}}
    x(t)\cos(2\pi kf_0t)\,dt,
    \quad k>0,
    \nonumber\\
    b_k
    &=2f_0\int_{-\frac{1}{2f_0}}^{\frac{1}{2f_0}}
    x(t)\sin(2\pi kf_0t)\,dt,
    \quad k>0.
\end{align}

% Q2.5
\item Find $a_k$ and $b_k$ for
\eqref{eq:10-orig-diff-def}.

\solution Since $x(t)$ is an even function,

\begin{align}
    x(-t)=x(t).
\end{align}

From the complex Fourier-series representation,

\begin{align}
    x(-t)
    &=\sum_{k=-\infty}^{\infty}
    c_ke^{-\jj2\pi kf_0t}\\
    &=\sum_{k=-\infty}^{\infty}
    c_{-k}e^{\jj2\pi kf_0t}.
\end{align}

Comparing this with

\begin{align}
    x(t)=\sum_{k=-\infty}^{\infty}
    c_ke^{\jj2\pi kf_0t},
\end{align}

we obtain

\begin{align}
    c_k=c_{-k}.
\end{align}

Therefore, from \eqref{eq:ak} and \eqref{eq:bk},

\begin{align}
    a_k
    &=
    \begin{cases}
        \dfrac{2A_0}{\pi}, & k=0,\\[6pt]
        \dfrac{4A_0}{\pi(1-k^2)}, &
        k>0,\ k\text{ even},\\[6pt]
        0, & k>0,\ k\text{ odd},
    \end{cases}
    \label{eq:ak-xt}\\
    b_k&=0,\qquad k>0.
    \label{eq:bk-xt}
\end{align}

For $A_0=12$, these become

\begin{align}
    a_0&=\frac{24}{\pi},\\
    a_k&=
    \begin{cases}
        \dfrac{48}{\pi(1-k^2)},&
        k>0,\ k\text{ even},\\[6pt]
        0,&k>0,\ k\text{ odd},
    \end{cases}\\
    b_k&=0.
\end{align}

% Q2.6
\item Verify \eqref{eq:one-Z-real} using Python.

\solution The Python code \texttt{codes/2\_6.py} verifies
\eqref{eq:one-Z-real} by comparing the original signal
with its real Fourier-series reconstruction, as shown in
Fig.~\ref{fig:ver-real}.

\begin{figure}[!ht]
    \centering
    \maybeimage{codes/2_6.png}
    \caption{Verification of the real Fourier series.}
    \label{fig:ver-real}
\end{figure}

\end{enumerate}
\section{Fourier Transform}

We use the Fourier-transform convention
\begin{align}
    G(f)=\int_{-\infty}^{\infty}g(t)e^{-\jj2\pi ft}\,dt,
    \label{eq:ft-definition}
\end{align}
with inverse transform
\begin{align}
    g(t)=\int_{-\infty}^{\infty}G(f)e^{\jj2\pi ft}\,df.
\end{align}

\begin{enumerate}[label=\thesection.\arabic*,ref=\thesection.\theenumi]

% Q3.1
\item The unit impulse satisfies
\begin{align}
    \delta(t)&=0,\quad t\neq 0,\\
    \int_{-\infty}^{\infty}\delta(t)\,dt&=1.
\end{align}

\solution The Dirac delta is zero everywhere except at the origin, but its total integral is unity. More generally, it has the sifting property
\begin{align}
    \int_{-\infty}^{\infty}g(t)\delta(t-t_0)\,dt=g(t_0).
\end{align}
This property allows an impulse to select the value of a function at its location.

% Q3.2
\item State the Fourier-transform pair for $g(t)$.

\solution By definition, the Fourier transform of $g(t)$ is
\begin{align}
    G(f)=\int_{-\infty}^{\infty}g(t)e^{-\jj2\pi ft}\,dt.
\end{align}
The corresponding inverse transform is
\begin{align}
    g(t)=\int_{-\infty}^{\infty}G(f)e^{\jj2\pi ft}\,df.
\end{align}
Hence, we write $g(t)\system{F}G(f)$.

% Q3.3
\item Show that
\begin{align}
    g(t-t_0)&\system{F}G(f)e^{-\jj2\pi ft_0}.
\end{align}

\solution Let $y(t)=g(t-t_0)$. Its Fourier transform is
\begin{align}
    Y(f)&=\int_{-\infty}^{\infty}g(t-t_0)e^{-\jj2\pi ft}\,dt.
\end{align}
Using the substitution $u=t-t_0$, so that $t=u+t_0$ and $dt=du$,
\begin{align}
    Y(f)&=\int_{-\infty}^{\infty}g(u)e^{-\jj2\pi f(u+t_0)}\,du\\
    &=e^{-\jj2\pi ft_0}\int_{-\infty}^{\infty}g(u)e^{-\jj2\pi fu}\,du\\
    &=G(f)e^{-\jj2\pi ft_0}.
\end{align}
Therefore,
\begin{align}
    \boxed{g(t-t_0)\system{F}G(f)e^{-\jj2\pi ft_0}}.
\end{align}

% Q3.4
\item Show that
\begin{align}
    G(t)&\system{F}g(-f).
\end{align}

\solution From the inverse Fourier-transform relation,
\begin{align}
    g(-f)=\int_{-\infty}^{\infty}G(\nu)e^{-\jj2\pi\nu f}\,d\nu.
\end{align}
This is precisely the Fourier transform of the function $G(t)$, evaluated at frequency $f$. Thus, by the duality property,
\begin{align}
    \boxed{G(t)\system{F}g(-f)}.
\end{align}

% Q3.5
\item Find the Fourier transform of $\delta(t)$.

\solution Using the definition,
\begin{align}
    \mathcal{F}\{\delta(t)\}
    &=\int_{-\infty}^{\infty}\delta(t)e^{-\jj2\pi ft}\,dt.
\end{align}
By the sifting property, the integrand is evaluated at $t=0$:
\begin{align}
    \boxed{\delta(t)\system{F}1}.
\end{align}

% Q3.6
\item Find the Fourier transform of $e^{-\jj2\pi f_0t}$.

\solution The standard frequency-shift impulse pair is
\begin{align}
    e^{\jj2\pi f_0t}&\system{F}\delta(f-f_0).
\end{align}
Replacing $f_0$ by $-f_0$ gives
\begin{align}
    \boxed{e^{-\jj2\pi f_0t}\system{F}\delta(f+f_0)}.
\end{align}

% Q3.7
\item Find the Fourier transform of $\cos(2\pi f_0t)$.

\solution Using Euler's identity,
\begin{align}
    \cos(2\pi f_0t)=\frac{1}{2}\left(e^{\jj2\pi f_0t}+e^{-\jj2\pi f_0t}\right).
\end{align}
By linearity and the exponential transform pairs,
\begin{align}
    \boxed{\cos(2\pi f_0t)\system{F}
    \frac{1}{2}\left[\delta(f-f_0)+\delta(f+f_0)\right]}.
\end{align}

% Q3.8
\item Find the Fourier transform of $x(t)$ and plot it. Verify using Python.

\solution The signal is periodic with fundamental frequency $2f_0$. Its complex Fourier-series representation can therefore be written as
\begin{align}
    x(t)=\sum_{m=-\infty}^{\infty}d_m e^{\jj2\pi(2mf_0)t},
\end{align}
where, from the coefficients obtained in Section 2,
\begin{align}
    d_m=c_{2m}=\frac{2A_0}{\pi(1-4m^2)}.
\end{align}
The Fourier transform of a periodic signal is an impulse train in the frequency domain. Hence,
\begin{align}
    X(f)&=\sum_{m=-\infty}^{\infty}d_m\delta(f-2mf_0)\\
    &=\boxed{\sum_{m=-\infty}^{\infty}
    \frac{2A_0}{\pi(1-4m^2)}\delta(f-2mf_0)}.
\end{align}
For $A_0=12$ and $f_0=50$ Hz,
\begin{align}
    X(f)=\sum_{m=-\infty}^{\infty}
    \frac{24}{\pi(1-4m^2)}\delta(f-100m).
\end{align}
Python code plots the first few spectral lines and compares the analytical Fourier-series reconstruction with the original signal.


\begin{figure}[!htp]
    \centering
    \maybeimage{codes/3_8.png}
    \caption{Comparison with the Fourier reconstruction.}
    \label{fig2}
\end{figure}



% Q3.9
\item Show that $\rect{t}\system{F}\sinc{f}$. Verify using Python.

\solution Using the normalized definitions
\begin{align}
    \rect{t}&=\begin{cases}1,&|t|<\frac{1}{2},\\0,&|t|>\frac{1}{2},\end{cases}\\
    \sinc{f}&=\frac{\sin(\pi f)}{\pi f},\qquad \sinc{0}=1,
\end{align}
the Fourier transform is
\begin{align}
    R(f)&=\int_{-\infty}^{\infty}\rect{t}e^{-\jj2\pi ft}\,dt\\
    &=\int_{-1/2}^{1/2}e^{-\jj2\pi ft}\,dt\\
    &=\left[\frac{e^{-\jj2\pi ft}}{-\jj2\pi f}\right]_{-1/2}^{1/2}\\
    &=\frac{e^{-\jj\pi f}-e^{\jj\pi f}}{-\jj2\pi f}\\
    &=\frac{\sin(\pi f)}{\pi f}.
\end{align}
Therefore,
\begin{align}
    \boxed{\rect{t}\system{F}\sinc{f}}.
\end{align}
The following code numerically evaluates the transform and compares it with the analytical sinc function.


\begin{figure}[!htp]
    \centering
    \maybeimage{codes/3_9.png}
    \caption{Comparison with the Fourier transform.}
    \label{fig3}
\end{figure}

% Q3.10
\item Find the Fourier transform of $\sinc{t}$. Verify using Python.

\solution By the duality property, since
\begin{align}
    \rect{t}&\system{F}\sinc{f},
\end{align}
we obtain
\begin{align}
    \sinc{t}&\system{F}\rect{-f}.
\end{align}
The rectangular function is even, so $\rect{-f}=\rect{f}$. Hence,
\begin{align}
    \boxed{\sinc{t}\system{F}\rect{f}}.
\end{align}
The following code numerically evaluates the transform over a finite time interval. Increasing the integration range and sampling density improves the approximation.
At the discontinuities $f=\pm\frac{1}{2}$, the Fourier-transform value is conventionally taken as the midpoint value $1/2$.

\begin{figure}[H]
    \centering
    \maybeimage{codes/3_10.png}
    \caption{Numerical verification of the Fourier transform of $\sinc{t}$.}
    \label{fig:fourier-sinc-verification}
\end{figure}

\end{enumerate}

\section{Filter}
\begin{enumerate}[label=\thesection.\arabic*
,ref=\thesection.\theenumi]
\item Find $H(f)$ which transforms $x(t)$ to DC 5V.
\item Find $h(t)$.
\item Verify your result using  through convolution.
\end{enumerate}
\section{Filter Design}
\begin{enumerate}[label=\thesection.\arabic*
,ref=\thesection.\theenumi]
\item Design a Butterworth filter for $H(f)$.
\item Design a Chebyschev filter for $H(f)$.
\item Design a circuit for your Butterworth filter.
\item Design a circuit for your Chebyschev filter.
\end{enumerate}




%------------------------------------------------

%------------------------------------------------
% The Fourier analysis above is intentionally two-column; the charger discussion is linear.
\onecolumn
\section{Circuit: From AC Input to Smoothed DC}
%------------------------------------------------

\subsection*{Purpose and signal path}
The circuit discussion follows the energy-conversion path up to the
smoothing capacitor. The stages have different jobs: the transformer
provides voltage reduction and isolation, the bridge rectifier makes
the load polarity unidirectional, and the capacitor reduces the
remaining periodic variation. The regulator and USB output are
discussed in the following sections.

\begin{itemize}[leftmargin=1.5em]
    \item \textbf{Input stage:} a 50\,Hz AC source is the reference input.
    \item \textbf{Voltage conversion and isolation:} a step-down
    transformer provides a lower AC voltage and galvanic isolation.
    \item \textbf{Rectification:} a four-diode bridge converts both
    half-cycles into positive output pulses.
    \item \textbf{Smoothing:} a reservoir capacitor stores energy near
    each peak and supplies the load between peaks.
    \item \textbf{Regulation:} a separate regulator produces a more
    stable 5\,V rail.
    \item \textbf{USB interface:} the port presents the correct
    electrical and charging-identification behaviour to a device.
\end{itemize}

\subsection*{Input quantities and notation}
The calculations below use the measured secondary voltage from the
input-stage analysis. This is kept distinct from the separate
12--0--12\,V transformer example in the reference circuit report.

\begin{table}[!t]
\centering
\caption{Input-stage values used in this report}
\label{tab:input-values}
\small
\begin{tabular}{|p{0.48\columnwidth}|c|c|}
\hline
\textbf{Quantity} & \textbf{Symbol} & \textbf{Value}\\ \hline
Primary RMS voltage & $V_p$ & $220\,\mathrm{V}$\\ \hline
Mains frequency & $f$ & $50\,\mathrm{Hz}$\\ \hline
Secondary RMS voltage (measured) & $V_s$ & $10.53\,\mathrm{V}$\\ \hline
One silicon diode drop (assumed) & $V_D$ & $0.7\,\mathrm{V}$\\ \hline
Bridge output ripple frequency & $f_r$ & $100\,\mathrm{Hz}$\\ \hline
Reservoir capacitor in main draft & $C$ & $2200\,\mu\mathrm{F}$\\ \hline
\end{tabular}
\end{table}

The symbols used throughout are:
\begin{itemize}[leftmargin=1.5em]
    \item $V_{\mathrm{rms}}$: root-mean-square voltage;
    \item $V_{\mathrm{pk}}$: peak voltage;
    \item $V_{\mathrm{avg}}$: average value over a cycle;
    \item $V_{r(pp)}$: peak-to-peak ripple voltage;
    \item $I_L$: load current; and
    \item $f_r$: frequency of the rectified ripple.
\end{itemize}

\subsection*{Mains waveform: RMS and peak values}
A sinusoidal voltage is commonly specified by its RMS value because
that value describes its equivalent heating effect in a resistive
load. For a pure sinusoid, the peak value is larger than the RMS value
by a factor of $\sqrt{2}$:
\begin{equation}
    \boxed{V_{\mathrm{pk}}=\sqrt{2}\,V_{\mathrm{rms}}.}
    \label{eq:rms-peak}
\end{equation}
For the nominal input stated in the reference calculation,
\begin{align}
    V_{p,\mathrm{pk}}
    &=220\sqrt{2}\\
    &\approx 311.1\,\mathrm{V}.
\end{align}
This value is included to explain the notation and the transformer
ratio; it is not an instruction to handle or measure mains voltage.

\textbf{Intuition.} RMS voltage is similar to an average
\emph{effectiveness} rating rather than the highest instantaneous
value. A 220\,V RMS sine wave does not remain at 220\,V; it swings
between positive and negative peaks of approximately 311\,V.

\subsection*{Step-down transformer}
A transformer transfers energy between windings through a changing
magnetic field. In an ideal transformer, the voltage ratio follows
the turns ratio:
\begin{equation}
    \boxed{\frac{V_s}{V_p}=\frac{N_s}{N_p}.}
    \label{eq:transformer-ratio}
\end{equation}
Here, $N_p$ and $N_s$ denote the primary and secondary turns. Using
the measured secondary RMS value:
\begin{align}
    \frac{N_p}{N_s}
    &=\frac{V_p}{V_s}
      =\frac{220}{10.53}\\
    &\approx 20.9:1.
\end{align}
The corresponding secondary peak is
\begin{equation}
    V_{s,\mathrm{pk}}=\sqrt{2}(10.53)
    \approx 14.89\,\mathrm{V}.
    \label{eq:secondary-peak}
\end{equation}

The transformer performs two conceptually separate functions:
\begin{itemize}[leftmargin=1.5em]
    \item \textbf{Voltage conversion:} the secondary voltage is lower
    than the primary voltage.
    \item \textbf{Galvanic isolation:} energy is transferred
    magnetically, without a direct conductive connection between
    primary and secondary windings in an appropriately designed
    isolation transformer.
\end{itemize}

\textbf{Contemporary analogy.} A transformer is like a contactless
gearbox for AC energy: it changes the voltage/current relationship
through magnetic coupling, much as gears trade speed for torque.
Unlike a simple resistor divider, it can transfer useful power
without intentionally burning the voltage difference as heat.

The reference main-circuit report describes a centre-tapped
12--0--12\,V transformer and reports an approximately 12\,V RMS
secondary waveform, corresponding to about 18\,V peak. That is a
different stated transformer condition from the measured
10.53\,V RMS value used in this report. These values must not be
combined in one calculation. The actual transformer label, winding
connection, load condition, and measured secondary voltage must be
confirmed before any design is finalised.

\subsection*{Full-wave bridge rectifier}
A bridge rectifier uses four diodes arranged so that the current
through the load has the same direction during both half-cycles of
the AC input. During one half-cycle, one diagonal pair conducts;
during the next half-cycle, the other pair conducts. At any instant,
the conducting path contains two forward-biased diodes.

\begin{itemize}[leftmargin=1.5em]
    \item Positive input half-cycle: the first diode pair carries
    current through the load.
    \item Negative input half-cycle: the second pair carries current
    through the load in the same load direction.
    \item Result: the negative half-cycle is flipped upward, producing
    a full-wave rectified waveform.
\end{itemize}

\textbf{Contemporary analogy.} Think of a one-way street system:
although the source alternates direction, the bridge routes both
directions of travel so that traffic through the load always moves
the same way.

For a silicon bridge, a first-order estimate uses two diode drops:
\begin{equation}
    \boxed{V_{o,\mathrm{pk}}\approx V_{s,\mathrm{pk}}-2V_D.}
    \label{eq:bridge-peak}
\end{equation}
With $V_D=0.7\,\mathrm{V}$ and the measured secondary peak:
\begin{align}
    V_{o,\mathrm{pk}}
    &=14.89-2(0.7)\\
    &\approx 13.49\,\mathrm{V}.
\end{align}
The ideal full-wave rectified sine has average value
\begin{equation}
    \boxed{V_{o,\mathrm{avg}}=\frac{2V_{o,\mathrm{pk}}}{\pi}}
    \approx 8.59\,\mathrm{V},
    \label{eq:rect-average}
\end{equation}
and RMS value
\begin{equation}
    \boxed{V_{o,\mathrm{rms}}=\frac{V_{o,\mathrm{pk}}}{\sqrt{2}}}
    \approx 9.54\,\mathrm{V}.
    \label{eq:rect-rms}
\end{equation}
These are idealised waveform calculations. Real diode forward
voltage depends on current and temperature, and transformer
regulation and winding resistance also affect the measured result.

Since both input half-cycles become positive pulses, the pulse
frequency doubles:
\begin{equation}
    \boxed{f_r=2f=100\,\mathrm{Hz},\qquad
    T_r=\frac{1}{f_r}=10\,\mathrm{ms}.}
    \label{eq:ripple-frequency}
\end{equation}
The bridge output is therefore \emph{pulsating DC}, not yet a steady
5\,V supply. Its voltage falls substantially between successive
peaks when a load is connected.

\textbf{Diode selection considerations.} The reference summary
compares the 1N4007 (nominal 1\,A average forward-current class and
1000\,V repetitive peak reverse-voltage rating) with the 1N5408
(3\,A average forward-current class). The correct diode or packaged
bridge must be selected using the actual load current, surge current,
reverse-voltage stress, temperature, and datasheet conditions.
A high reverse-voltage rating alone does not establish suitability
for the required current or thermal conditions.

\begin{table}[!t]
\centering
\caption{Calculated waveform values at the bridge output}
\label{tab:bridge-results}
\small
\begin{tabular}{|p{0.48\columnwidth}|c|c|}
\hline
\textbf{Quantity} & \textbf{Symbol} & \textbf{Value}\\ \hline
Secondary RMS voltage & $V_s$ & $10.53\,\mathrm{V}$\\ \hline
Secondary peak voltage & $V_{s,\mathrm{pk}}$ & $14.89\,\mathrm{V}$\\ \hline
Bridge peak voltage & $V_{o,\mathrm{pk}}$ & $13.49\,\mathrm{V}$\\ \hline
Bridge average voltage & $V_{o,\mathrm{avg}}$ & $8.59\,\mathrm{V}$\\ \hline
Bridge RMS voltage & $V_{o,\mathrm{rms}}$ & $9.54\,\mathrm{V}$\\ \hline
Ripple frequency & $f_r$ & $100\,\mathrm{Hz}$\\ \hline
Ripple period & $T_r$ & $10\,\mathrm{ms}$\\ \hline
\end{tabular}
\end{table}

\subsection*{Smoothing capacitor: reservoir action}
A capacitor connected across the rectifier output acts as an energy
reservoir. Near each rectified peak, the rectifier voltage exceeds
the capacitor voltage and the capacitor is replenished. As the
rectifier voltage falls below the capacitor voltage, the diodes stop
conducting and the capacitor supplies the load until the next peak.

\begin{itemize}[leftmargin=1.5em]
    \item \textbf{Charging interval:} the capacitor voltage rises
    towards the rectifier peak.
    \item \textbf{Discharging interval:} the capacitor releases stored
    energy into the load between peaks.
    \item \textbf{Ripple:} the repeated charge/discharge cycle creates
    a small periodic variation superimposed on the DC level.
    \item \textbf{Larger capacitance:} for the same load current and
    ripple frequency, a larger capacitor generally reduces the
    voltage drop between peaks.
\end{itemize}

\textbf{Contemporary analogy.} The capacitor behaves like a small
water tank placed after an intermittent pump. The rectifier refills
the tank at each pulse; the load draws water continuously. A larger
tank reduces the change in water level between refills. The analogy
describes energy buffering, not a literal equivalence between
voltage and water pressure.

For a resistive load $R_L$, the ideal discharge between charging
peaks is described by
\begin{equation}
    \boxed{v_C(t)=V_{\max}
    e^{-(t-t_0)/(R_LC)}.}
    \label{eq:cap-discharge}
\end{equation}
The associated time constant is
\begin{equation}
    \boxed{\tau=R_LC.}
    \label{eq:rc-time}
\end{equation}
A larger $\tau$ means the capacitor voltage decays more slowly
during the interval between rectifier peaks.

For a full-wave rectifier and relatively small ripple, the
peak-to-peak ripple can be estimated by
\begin{equation}
    \boxed{V_{r(pp)}\approx\frac{I_L}{f_rC}.}
    \label{eq:ripple-estimate}
\end{equation}
For $f_r=100\,\mathrm{Hz}$ and $C=2200\,\mu\mathrm{F}$:
\begin{align}
    V_{r(pp)}
    &\approx \frac{I_L}{100(2200\times10^{-6})}\\
    &\approx 4.545 I_L\ \mathrm{V},
\end{align}
where $I_L$ is expressed in amperes. The approximate average
capacitor voltage is
\begin{equation}
    \boxed{V_{\mathrm{DC}}\approx
    V_{\max}-\frac{V_{r(pp)}}{2}.}
    \label{eq:cap-average}
\end{equation}
With $V_{\max}\approx13.49\,\mathrm{V}$, this becomes
\begin{equation}
    V_{\mathrm{DC}}\approx13.49-2.273I_L\ \mathrm{V}.
\end{equation}
This is an estimate, not a measured output. It assumes small ripple
and nearly constant load current. The load current, capacitor
tolerance, equivalent series resistance (ESR), transformer
impedance, diode conduction interval, and temperature all affect
the actual ripple.

The reference circuit report mentions a $100\,\mu\mathrm{F}$
capacitor as its filter element and reports an approximately
18\,V DC level in its particular simulation/measurement. The
present main draft uses $2200\,\mu\mathrm{F}$ in its calculations.
Those are not interchangeable experimental records. The capacitor
value must be chosen from the intended load current and allowed
ripple, and its voltage and ripple-current ratings must be suitable
for the actual circuit. The reported 18\,V result belongs to the
reference's 12\,V RMS / 18\,V peak example, not to the
10.53\,V RMS calculation above.

\textbf{Important limitation.} A reservoir capacitor does not
produce a precision DC output. It reduces the low-frequency ripple
but does not guarantee a fixed voltage as the load changes. It also
does not replace the regulator, current limiting, thermal design,
or USB source-identification circuitry.

%------------------------------------------------
\section{Regulator and Buck Converter}
%------------------------------------------------
The regulator stage converts the smoothed, load-dependent DC rail
into a controlled output. The two devices considered here are
alternatives, not stages that must be connected in series. Both
can provide a nominal 5\,V rail, but their energy-conversion
mechanisms and practical limitations are different.

\subsection{IC 7805 linear regulator}
The 7805 is a fixed positive linear regulator. Its internal reference
and error amplifier monitor the output and control a series pass
transistor. If the output tends to fall under increased load, the
control circuit increases conduction; if the output tends to rise,
it reduces conduction. This feedback action keeps the output near
5\,V while the device remains within its operating and thermal
limits.

\begin{itemize}[leftmargin=1.5em]
    \item Fixed nominal output: $V_{\mathrm{OUT}}\approx5\,\mathrm{V}$.
    \item Requires input headroom above 5\,V; a useful first estimate
    is $V_{\mathrm{IN}}\gtrsim V_{\mathrm{OUT}}+V_{\mathrm{dropout}}$.
    \item The TI LM340/LM7805 family documentation gives about
    2\,V dropout at a stated 1\,A test condition. The exact value
    depends on device, load, and temperature.
    \item The excess voltage is mainly dissipated as heat in the
    pass element.
\end{itemize}

The approximate power dissipation is
\begin{equation}
    \boxed{P_D\approx(V_{\mathrm{IN}}-V_{\mathrm{OUT}})I_{\mathrm{OUT}}.}
    \label{eq:7805-loss}
\end{equation}
For an illustrative 12\,V input, 5\,V output, and 0.5\,A load:
\begin{equation}
    P_D\approx(12-5)(0.5)=3.5\,\mathrm{W}.
\end{equation}
At 1\,A under the same voltage conditions, the estimate rises to
7\,W. Such dissipation can make a small package very hot unless
the thermal path is designed appropriately.

\textbf{Thermal intuition.} A 7805 is like a pressure-reducing valve
that discards the excess pressure. The output is controlled, but
the unused voltage difference is not recovered; electrically, much
of that difference becomes heat.

A simplified junction-temperature estimate is
\begin{equation}
    \boxed{T_J\approx T_A+P_D R_{\theta JA},}
    \label{eq:thermal}
\end{equation}
where $T_A$ is ambient temperature and $R_{\theta JA}$ is the
effective junction-to-ambient thermal resistance of the assembled
device. A heatsink or larger copper area can improve heat removal,
but the actual thermal resistance depends on mounting and airflow.

The 7805 family includes current limiting, thermal shutdown, and
safe-operating-area protection. These are protective functions,
not permission to operate continuously at the protection threshold.
A family capability above 1.5\,A is conditional on device variant
and thermal design; it is not a guaranteed continuous current for
every low-cost module.

\textbf{Advantages}
\begin{itemize}[leftmargin=1.5em]
    \item Simple circuit and straightforward fixed 5\,V output.
    \item No switching-frequency ripple generated by the regulator.
    \item Useful when the input-output voltage difference and load
    current are modest.
\end{itemize}
\textbf{Limitations}
\begin{itemize}[leftmargin=1.5em]
    \item Efficiency becomes poor when the input is much higher
    than 5\,V.
    \item Dissipation increases with both voltage difference and
    output current.
    \item Input must remain above the required headroom even at
    ripple valleys and under load.
\end{itemize}

\subsection{LM2596 buck converter}
The LM2596 is a fixed-frequency step-down switching regulator.
The family includes fixed-output versions and an adjustable version.
The TI datasheet specifies a 4.5--40\,V operating input range,
a 3\,A output-current capability under stated conditions, and a
150\,kHz oscillator. The ratings of a complete module also depend
on its inductor, diode, capacitors, PCB layout, and cooling.

A buck converter does not continuously burn off the excess input
voltage. Instead, an internal switch transfers energy in pulses to
an inductor. The inductor stores energy in its magnetic field and
continues delivering current when the switch is off. A catch diode
provides the current path during the off interval, while the output
capacitor reduces the resulting voltage variation.

\begin{itemize}[leftmargin=1.5em]
    \item \textbf{Switch ON:} the internal switch connects the input
    to the inductor path; inductor current rises and energy is stored.
    \item \textbf{Switch OFF:} the switch opens; inductor current
    continues through the catch diode and supplies the load.
    \item \textbf{Output filtering:} the capacitor supplies or absorbs
    the difference between inductor current and load current.
    \item \textbf{Feedback:} the controller senses the output and
    adjusts switching behaviour to keep the average output near
    the set value.
\end{itemize}

\textbf{Contemporary analogy.} The inductor acts like a flywheel:
it resists abrupt changes in current and keeps energy moving during
the switch-off interval. The capacitor acts like a buffer tank for
voltage. Together they turn a train of switching pulses into a
much steadier output.

For an ideal buck converter operating in continuous-conduction
mode, the average output is
\begin{equation}
    \boxed{V_{\mathrm{OUT}}=D V_{\mathrm{IN}},\qquad
    D=\frac{V_{\mathrm{OUT}}}{V_{\mathrm{IN}}},}
    \label{eq:buck-duty}
\end{equation}
where $D$ is the duty ratio. For an ideal 12\,V to 5\,V conversion:
\begin{equation}
    D\approx\frac{5}{12}=0.417=41.7\%.
\end{equation}
Real output differs because of switch and diode drops, winding
resistance, ripple, control limits, and load conditions.

The efficiency is
\begin{equation}
    \boxed{\eta=\frac{P_{\mathrm{OUT}}}{P_{\mathrm{IN}}}.}
    \label{eq:efficiency}
\end{equation}
For illustration only, a 5\,V, 1\,A output is 5\,W. At an assumed
85\% efficiency, the input power would be approximately
$5/0.85=5.88$\,W and the loss approximately 0.88\,W. This is an
example, not a guaranteed efficiency figure for every LM2596
module.

\textbf{Advantages}
\begin{itemize}[leftmargin=1.5em]
    \item Usually more efficient than a linear regulator for a
    substantial step-down ratio.
    \item Lower heat dissipation for comparable output power in
    many operating conditions.
    \item Fixed and adjustable output variants are available.
\end{itemize}
\textbf{Limitations}
\begin{itemize}[leftmargin=1.5em]
    \item Requires an inductor, catch diode, and suitable input and
    output capacitors.
    \item Switching ripple and electromagnetic interference (EMI)
    must be considered.
    \item The 3\,A IC rating does not mean every inexpensive module
    can deliver 3\,A continuously.
    \item Inductor saturation, diode current, capacitor ripple
    rating, PCB copper, and temperature constrain usable output.
\end{itemize}

\begin{table}[!t]
\centering
\caption{Practical comparison of the two 5\,V regulator options}
\label{tab:reg-comparison}
\small
\begin{tabular}{|p{0.25\columnwidth}|p{0.31\columnwidth}|p{0.31\columnwidth}|}
\hline
\textbf{Feature} & \textbf{7805} & \textbf{LM2596}\\ \hline
Method & Linear & Switching buck\\ \hline
Output & Fixed 5\,V & Fixed or adjustable\\ \hline
Main loss & Voltage drop as heat & Switching and component losses\\ \hline
External parts & Bypass capacitors; thermal path & Inductor, diode, capacitors\\ \hline
Ripple character & No switching ripple & Switching ripple present\\ \hline
Best fit & Small drop, modest current & Larger voltage reduction\\ \hline
Main design check & Dissipation and headroom & Current, ripple, layout, heat\\ \hline
\end{tabular}
\end{table}

\textbf{Selection for this project.} If the smoothed input is
substantially above 5\,V and the USB load current is meaningful,
the LM2596 is generally the more practical candidate because it
avoids the large linear loss predicted by
\eqref{eq:7805-loss}. A 7805 remains useful as a simple educational
reference or for a low-current rail when its headroom and thermal
conditions are verified. The regulator must be selected from the
actual DC input range, not from the transformer's nominal label
alone.

%------------------------------------------------
\section{USB Output Port}
%------------------------------------------------
A USB connector is not merely a pair of 5\,V and ground contacts.
It is an interface with power, return, data or identification
contacts, mechanical requirements, and rules that tell a connected
device how much current it may draw. A regulator can hold a 5\,V
rail, but it does not automatically make the port a compliant
charger.

\textbf{Main functions of the USB output}
\begin{itemize}[leftmargin=1.5em]
    \item \textbf{VBUS:} carries the positive supply, nominally 5\,V
    for the basic USB charging modes discussed here.
    \item \textbf{GND:} provides the return path for current.
    \item \textbf{D+ and D-:} USB 2.0 data lines; in legacy charging
    arrangements, their electrical signature also helps a device
    identify the source type.
    \item \textbf{CC pins on USB Type-C:} advertise source current
    and establish connection orientation/role information.
    \item \textbf{USB-PD communication:} negotiates supported
    voltage/current profiles when a suitable Power Delivery
    controller and power path are present.
\end{itemize}

\textbf{Contemporary analogy.} VBUS is the delivery lane carrying
energy. The identification pins are the signboard at the entrance
that tells the device what kind of lane it has entered and what
current allowance is available. A powerful supply without the
correct signboard may still be treated as a lower-capability source.

\textbf{Voltage is not the same as current.} The regulator establishes
the output voltage. The connected device draws current according
to its operating state and the source capability it detects. The
charger must be able to maintain its voltage at that load, but it
should not force a fixed current into every device. Current
advertisement, current limiting, cable capability, and protection
are separate design requirements.

\textbf{Legacy USB downstream ports.}
A Standard Downstream Port (SDP) is a normal host/data port. The
USB 2.0 reference limit is up to 500\,mA after the applicable
enumeration/configuration process; a USB 3.x SDP is up to 900\,mA
under its specification. These are not blanket promises that any
device can immediately draw those currents from any USB socket.
Host rules and device configuration apply.

\textbf{Battery Charging Specification 1.2.}
BC1.2 distinguishes three common port types:
\begin{itemize}[leftmargin=1.5em]
    \item \textbf{SDP (Standard Downstream Port):} supports USB data
    and follows the applicable USB host current rules.
    \item \textbf{CDP (Charging Downstream Port):} supports data and
    provides charging capability, with up to 1.5\,A under the
    standard's conditions.
    \item \textbf{DCP (Dedicated Charging Port):} intended for
    charging rather than USB data; the D+ and D- lines have the
    BC1.2 DCP signature, commonly implemented as a short between
    the two lines. A compatible device can identify the port and
    draw up to 1.5\,A within the specification.
\end{itemize}
A device must perform the relevant detection. Merely shorting pins
or applying a voltage without following the applicable standard
does not make a universal or compliant charging port.

\begin{table}[!t]
\centering
\caption{Common USB charging port categories}
\label{tab:usb-port-types}
\small
\begin{tabular}{|p{0.22\columnwidth}|p{0.35\columnwidth}|p{0.28\columnwidth}|}
\hline
\textbf{Port type} & \textbf{Data support} & \textbf{Reference current}\\ \hline
USB 2.0 SDP & Yes & Up to 500\,mA*\\ \hline
USB 3.x SDP & Yes & Up to 900\,mA*\\ \hline
BC1.2 CDP & Yes & Up to 1.5\,A\\ \hline
BC1.2 DCP & No data function & Up to 1.5\,A\\ \hline
\end{tabular}
\end{table}
\smallskip
\noindent\textit{*Host enumeration and configuration rules apply.}

\textbf{Legacy proprietary D+/D- signatures.}
Some older charging ecosystems used voltage levels on D+ and D-
to identify a source's current capability. The main draft's
reference table records example Apple-style signatures. These
values are legacy, proprietary arrangements and should not be
treated as universal USB rules.

\begin{table}[!t]
\centering
\caption{Legacy D+/D- voltage signatures recorded in the reference}
\label{tab:legacy-usb}
\small
\begin{tabular}{|c|c|c|}
\hline
\textbf{D+} & \textbf{D-} & \textbf{Reported current}\\ \hline
2.0\,V & 2.0\,V & Approx. 0.5\,A\\ \hline
2.0\,V & 2.7\,V & Approx. 1.0\,A\\ \hline
2.7\,V & 2.0\,V & Approx. 2.1\,A\\ \hline
2.7\,V & 2.7\,V & Approx. 2.4\,A\\ \hline
\end{tabular}
\end{table}
These are included to explain why older devices sometimes used
different charging signatures. They are not a recommended
implementation for a new universal charger. Device support varies,
and the applicable source documentation must be checked.

\textbf{USB Type-C current advertisement.}
Type-C uses Configuration Channel (CC) pins rather than relying
only on the legacy D+/D- signatures. A compliant 5\,V Type-C source
can advertise default USB current, 1.5\,A, or 3\,A through the
appropriate CC termination. The connected device reads that
advertisement and draws only within the advertised capability.
The cable and complete power path must also be rated for the
current.

\textbf{USB Power Delivery (USB-PD).}
USB-PD is a separate communication protocol. It allows a source
and sink to negotiate supported voltage and current profiles.
A USB-C connector connected to a fixed 5\,V regulator does not
automatically provide USB-PD. A suitable PD controller, compliant
power-path design, cable, and negotiation process are required
for higher-voltage profiles.

\begin{table}[!t]
\centering
\caption{What the regulator does and what the USB interface must do}
\label{tab:usb-responsibility}
\small
\begin{tabular}{|p{0.43\columnwidth}|p{0.43\columnwidth}|}
\hline
\textbf{Regulator / power path} & \textbf{USB interface}\\ \hline
Maintain the selected DC voltage & Identify the port type\\ \hline
Supply the required load current & Advertise or negotiate allowed current\\ \hline
Remain within thermal limits & Use suitable connector and cable\\ \hline
Provide appropriate electrical protection & Respect USB protocol requirements\\ \hline
\end{tabular}
\end{table}

\textbf{Output protection and verification.}
A practical low-voltage USB output needs a correctly rated current
limit, short-circuit protection, overvoltage protection where
appropriate, reverse-current considerations, suitable connector
and cable ratings, and a stable output under load. The regulator's
headline current rating alone does not establish the safe current
capability of the complete port. Verification should use an
appropriate electronic load and measurement equipment on an
isolated low-voltage output; no USB output should be tested while
the upstream mains side is exposed.

%------------------------------------------------
\section{Engineering Review and Design Checklist}
%------------------------------------------------
The preceding analysis shows that each stage has a distinct
requirement. A correct calculation at one stage does not prove that
the complete charger is suitable. The values must be checked as a
connected system under realistic load and temperature conditions.

\begin{enumerate}[leftmargin=1.6em]
    \item \textbf{Input definition:} record the actual secondary RMS
    voltage and frequency under the intended load. Do not mix the
    measured 10.53\,V value with the separate 12--0--12\,V reference
    example.
    \item \textbf{Rectifier:} verify the bridge current and reverse
    voltage ratings, surge capability, diode losses, and temperature.
    \item \textbf{Capacitor:} select capacitance from load current
    and allowable ripple; check voltage rating, polarity, ESR, and
    ripple-current rating.
    \item \textbf{Regulator:} check the minimum input at ripple
    valleys, maximum input at peaks, output current, efficiency,
    dissipation, and thermal behaviour.
    \item \textbf{USB port:} choose one supported charging standard
    and implement its identification/advertisement correctly.
    \item \textbf{System verification:} check output voltage, ripple,
    temperature, and protection behaviour at no load and at the
    intended low-voltage test loads.
\end{enumerate}

\textbf{Source-data reconciliation.}
The reference circuit report describes a centre-tapped 12--0--12\,V
transformer, a 100\,$\mu$F filter capacitor, and an approximately
18\,V peak/filtered level in its reported setup. The input-stage
summary supplied with this project instead uses a measured
10.53\,V RMS secondary, calculates 14.89\,V peak before the bridge,
and 13.49\,V peak after two diode drops. The main source file
specifies 2200\,$\mu$F. These are different data sets. This draft
retains the 10.53\,V measured case and 2200\,$\mu$F as the main
calculation, while explicitly identifying the 12--0--12\,V / 100
$\mu$F values as reference-report data requiring confirmation.

The reference circuit's reported waveform labels also call the
rectifier waveform a 12\,V, 50\,Hz rectified sinusoid in one place,
while the results list an 18\,V peak after transformer and
rectifier. A full-wave bridge supplied at 50\,Hz produces a
100\,Hz ripple waveform. The 12\,V and 18\,V descriptions are
therefore not sufficient by themselves to establish a consistent
measured waveform. The calculation in this report uses the explicit
RMS/peak equations and labels the resulting ripple frequency
correctly.

\section*{References}
%------------------------------------------------
\begin{enumerate}[leftmargin=1.5em]
    \item Texas Instruments, \emph{LM340/LM7805 Family Datasheet},
    Wide VIN 1.5-A Fixed Voltage Regulators, Rev. L.
    \url{https://www.ti.com/lit/ds/symlink/lm7800.pdf}
    \item Texas Instruments, \emph{LM2596 SIMPLE SWITCHER 150-kHz
    3-A Step-Down Voltage Regulator}, Datasheet, Rev. G.
    \url{https://www.ti.com/lit/ds/symlink/lm2596.pdf}
    \item USB Implementers Forum, \emph{USB 2.0 Specification}
    and compliance documentation.
    \url{https://www.usb.org/document-library/usb-20-specification}
    \item Microchip Technology, \emph{Battery Charging with USB:
    SDP, CDP and DCP Overview}.
    \url{https://developerhelp.microchip.com/xwiki/bin/view/applications/usb/power-delivery/battery-charging/}
    \item USB Battery Charging Specification 1.2.
    \url{https://www.usb.org/document-library/usb-20-specification}
\end{enumerate}

\end{document}

